\section{Выбор \latGCKgen\ с двенадцатью соседями из правила «один такт = одно ребро»}
\label{sec:fcc-choice}

В~\cref{sec:epsilon} установлено главное правило дискретного переноса: \emph{перемещаться можно только по рёбрам каузальной окрестности $N(x)$ --- один такт~$\htick$ равен одному ребру}. Если между двумя узлами нет ребра, однотактовый переход между ними невозможен. Посмотрим, что из этого правила с необходимостью следует для размерности пространства~$d$ и устройства самой оболочки~$N(x)$.

Пусть в $d$-мерном пространстве заданы $d$ единичных базисных векторов $\mathbf{e}_1,\dots,\mathbf{e}_d$.
\begin{enumerate}
\item \textbf{Откуда берётся уравнение Каталана $1+2^d=d^2$ и размерность $d=3$:}
  на планковском полу (\cref{sec:planck-floor,sec:descent-phase-disc}) неделимая элементарная ячейка~$\caPlanckCell$ не может быть заполнена дробно --- она либо свободна ($0$), либо занята ($1$). Поэтому внутри одного узла собственное состояние выбора складывается из $1$~бита занятости самой центральной ячейки (остаться на месте вне осей) и $2^d$ бинарных комбинаций участия (свободно/занято) вдоль $d$ единичных векторов $\mathbf{e}_1,\dots,\mathbf{e}_d$ --- всего $1+2^d$ состояний. С другой стороны, любой элементарный акт движения или поворота на решётке есть перенос занятости с одного единичного вектора $\mathbf{e}_i$ на другой единичный вектор $\mathbf{e}_j$; таких упорядоченных пар единичных векторов $(\mathbf{e}_i\to\mathbf{e}_j)$ существует ровно $d\times d = d^2$. Чтобы узел кодировал все элементарные переходы $(\mathbf{e}_i\to\mathbf{e}_j)$ своим собственным состоянием без внешней памяти, число состояний выбора обязано совпадать с числом пар переходов:
  \begin{equation}
    \label{eq:catalan-from-edges}
    1 + 2^d \;=\; d^2 \quad\Longleftrightarrow\quad d = 3 \qquad \bigl(1 + 2^3 = 9 = 3^2\bigr).
  \end{equation}
\item \textbf{Почему поворот обязан быть ребром $\pm\mathbf{e}_i\pm\mathbf{e}_j$ ($4+4+4=12$ рёбер кубооктаэдра):}
  разложим $d^2=9$ пар переходов $(\mathbf{e}_i\to\mathbf{e}_j)$ на $d=3$ продольные пары $(\mathbf{e}_i\to\mathbf{e}_i)$ вдоль одной и той же оси и $d(d-1)=6$ внедиагональных пар $(\mathbf{e}_i\to\mathbf{e}_j)$ ($i\ne j$), переводящих движение с оси~$i$ на другую ось~$j$:
  \[
    d^2 \;=\; d + d(d-1) \;=\; 3 + 6 \;=\; 9.
  \]
  Если один такт равен одному ребру, то \emph{поворот обязан быть ребром}. Простая кубическая сетка (\latPKshort, рёбра $\pm\mathbf{e}_i$, $|N|=6$) составлена только из продольных пар $(\mathbf{e}_i\to\mathbf{e}_i)$ и не содержит ни одного ребра, меняющего направление за один шаг. Чтобы поворот с $\mathbf{e}_i$ на $\mathbf{e}_j$ ($i\ne j$) осуществлялся за один такт, само ребро обязано иметь вид
  \begin{equation}
    \label{eq:turn-edge-cuboct}
    e \;=\; \pm\mathbf{e}_i \pm\mathbf{e}_j \qquad (1\le i < j \le d).
  \end{equation}
  В трёх измерениях ($d=3$) число независимых плоскостей вращения и наклона (число углов Эйлера $\dim\mathrm{SO}(d)=d(d-1)/2$) в точности совпадает с числом осей $d=3$: это три координатные плоскости $(xy,xz,yz)$. В каждой из трёх плоскостей лежат $4$ знаковые комбинации~\eqref{eq:turn-edge-cuboct}, что даёт
  \[
    |N(x)| \;=\; 4 + 4 + 4 \;=\; 2d(d-1) \;=\; 12
  \]
  поворотных рёбер строго одинаковой евклидовой длины $\|\pm\mathbf{e}_i\pm\mathbf{e}_j\|_2 = \sqrt{1^2+1^2} = \sqrt{2}$ (в единицах полушага куба), концы которых образуют \textbf{кубооктаэдр} гранецентрированной кубической решётки (\latGCKshort).
\end{enumerate}

\begin{definition}[Типы упаковки в трёх измерениях]
\label{def:packing-types}
Сравнивая носители, используем стандартную кристаллографическую терминологию
(без англоязычных аббревиатур SC, BCC, FCC, HCP в прозе):
\begin{center}
\small
\begin{tabular}{@{}p{0.52\linewidth}cc@{}}
\toprule
название & сокр. & $|N|$ \\
\midrule
\latPK & \latPKshort & $6$ \\
\latOCK & \latOCKshort & $8$ \\
\latGCK & \latGCKshort & $12$ \\
\latGPU & \latGPUshort & $12$ \\
\bottomrule
\end{tabular}
\end{center}
Носитель модели --- \latGCK\ $\carrierLattice\subset\mathbb{R}^3$.
В формулах иногда пишут $\mathrm{FCC}$ как условное обозначение этого типа Браве.
\end{definition}

\begin{theorem}[Носитель в $(3{+}1)$]
\label{thm:fcc-choice}
Пусть выполнены: замощение без дыр; единственная длина~$\caStepSpace$;
только светоподобные рёбра за один~$\htick$;
изотропия вакуума; один носитель на узел; физическое пространство трёхмерно.
Тогда $\carrierLattice$ --- \latGCK, $|N(x)|=12$ для всех узлов,
$\kappa_{\mathrm{link}}=1/12$.
\end{theorem}

\begin{proof}
Сначала отбираем упаковки с одной длиной ребра и без дыр, затем отсекаем вторую оболочку световым конусом.

\emph{Упаковка.}
\pusing{модели шара радиуса $\caStepSpace/2$ на узле}
Ячейку представляем шаром радиуса $\caStepSpace/2$: касание двух шаров есть ребро длины~$\caStepSpace$.
В трёх измерениях плотнейшая упаковка равных шаров --- гранецентрированная кубическая или гексагональная плотная (теорема Кеплера).
У каждого шара тогда ровно двенадцать соседей на расстоянии~$\caStepSpace$.
\latPK\ имеет шесть соседей, \latOCK --- восемь; обе реже и отпадают, раз требуется максимальная плотность при одной длине ребра
(\cref{fig:carrier-packing-sc,fig:carrier-packing-bcc,fig:carrier-packing-fcc}).

\begin{figure}[htbp]
  \centering
  \input{figures/tikz/carrier-packing-sc.tex}
  \caption{\latPK\ (\latPKshort): $|N|=6$.}
  \label{fig:carrier-packing-sc}
\end{figure}

\begin{figure}[htbp]
  \centering
  \input{figures/tikz/carrier-packing-bcc.tex}
  \caption{\latOCK\ (\latOCKshort): $|N|=8$.}
  \label{fig:carrier-packing-bcc}
\end{figure}

\begin{figure}[htbp]
  \centering
  \input{figures/tikz/carrier-packing-fcc.tex}
  \caption{\latGCK\ (\latGCKshort): $|N|=12$.}
  \label{fig:carrier-packing-fcc}
\end{figure}

Плотность и $|N|$ у \latGCKshort\ и \latGPUshort\ совпадают; их различает симметрия.
\pfrom{требования изотропии вакуума}
у \latGCKshort\ три равные кубические оси, у \latGPUshort --- выделенная ось слоёв.
\phence носитель --- \latGCK.

\emph{Световой конус.}
Двенадцать соседей узла в начале координат, при ребре кубической ячейки $\caStepSpace\sqrt{2}$, суть все перестановки
\[
  \bigl(\pm\caStepSpace/\sqrt{2},\ \pm\caStepSpace/\sqrt{2},\ 0\bigr).
\]
Расстояние до начала
\[
  \sqrt{2\cdot(\caStepSpace/\sqrt{2})^2}=\caStepSpace.
\]
Вторая оболочка --- перестановки $(\pm\caStepSpace\sqrt{2},0,0)$, расстояние $\caStepSpace\sqrt{2}$.
\pusing{одного такта $\caSignalSpeed\htick=\caStepSpace$}
интервал до узла второй оболочки
\[
  \caSignalSpeed^2\htick^2-(\caStepSpace\sqrt{2})^2=\caStepSpace^2-2\caStepSpace^2=-\caStepSpace^2<0.
\]
\pfollows
узел второй оболочки за этот такт недостижим
(\cref{fig:carrier-light-cone-2d,fig:carrier-light-cone-3d}).
Остаётся первая оболочка: $|N(x)|=12$ (\cref{fig:carrier-fcc-shell}).

\begin{figure}[htbp]
  \centering
  \input{figures/tikz/carrier-light-cone-2d.tex}
  \caption{Световой конус $\caSignalSpeed$ в срезе $(x,t)$: первая оболочка на конусе, вторая ($\caStepSpace\sqrt{2}$) --- вне.}
  \label{fig:carrier-light-cone-2d}
\end{figure}

\begin{figure}[htbp]
  \centering
  \input{figures/tikz/carrier-light-cone-3d.tex}
  \caption{Световой конус в объёме $(x,y,t)$.}
  \label{fig:carrier-light-cone-3d}
\end{figure}

\begin{figure}[htbp]
  \centering
  \input{figures/tikz/carrier-fcc-shell.tex}
  \caption{Плотнейшая \latGCK: $|N|=12$ касающихся сфер.}
  \label{fig:carrier-fcc-shell}
\end{figure}

\emph{Связь.}
\pfrom{изотропии микрошага на двенадцати рёбрах}
\[
  \kappa_{\mathrm{link}}=\frac{1}{|N|}=\frac{1}{12}.
\]
\end{proof}

\subsection{Однотактовое тело: почему кубооктаэдр}
\label{sec:one-tick-hull}

Кубооктаэдр в тексте появляется не как «выбранный из каталога Архимеда'' многогранник,
а как \emph{имя} пространственного следа одного микрошага при уже принятом носителе.

\begin{definition}[Однотактовая оболочка на $\Mlayer$]
\label{def:one-tick-hull}
Пусть $N(x)=\{y:|y-x|=\caStepSpace\}$ --- первая координационная сфера (\cref{thm:fcc-choice}).
При условии~\ref{ax:A1} за один~$\htick$ из~$x$ достижимы только точки~$N(x)$.
Пространственное тело одного такта --- выпуклая оболочка $\mathcal{K}(x)=\mathrm{conv}\,N(x)$.
\end{definition}

На альтернативных решётках, отсечённых в доказательстве~\ref{thm:fcc-choice}, то же определение
даёт другие тела --- и это показывает, что многогранник \emph{следует} за решёткой, а не подменяет её:
\begin{center}
\small
\begin{tabular}{@{}lll@{}}
\toprule
решётка & $|N|$ & $\mathcal{K}(x)=\mathrm{conv}\,N(x)$ \\
\midrule
\latPKshort & $6$ & октаэдр \\
\latOCKshort & $8$ & куб \\
\latGCKshort\ / \latGPUshort & $12$ & кубооктаэдр / призма с шестиугольным основанием \\
\bottomrule
\end{tabular}
\end{center}
\latGPU\ с той же плотностью, что \latGCK, оставляет выделенную ось слоёв;
\pfrom{изотропии вакуума в \cref{thm:fcc-choice}}
носитель --- \latGCK, а не \latGPU.

\begin{proposition}[Кубооктаэдр как оболочка \latGCKshort, три плоскости углов Эйлера и закон изотропии $O(|\mathbf{k}|^2)$ с анизотропией $O(|\mathbf{k}|^4)$]
\label{prop:cubocta-fcc}
Пусть $\carrierLattice$ --- \latGCK\ и $a=\caStepSpace$.
Тогда $N(0)$ --- все двенадцать поворотных векторов $(\pm a/\sqrt{2},\pm a/\sqrt{2},0)$ (в безразмерных полушагах куба: $O_2(3)=\{\pm\mathbf{e}_i\pm\mathbf{e}_j \mid 1\le i<j\le 3\}$),
и $\mathcal{K}(0)=\mathrm{conv}\,N(0)$ --- правильный кубооктаэдр с ребром~$a$,
шестью квадратными и восемью треугольными гранями (\cref{fig:carrier-cuboctahedron}), обладающий следующими свойствами:
\begin{enumerate}[label=(\roman*)]
\item \textbf{Три плоскости углов Эйлера и треугольные шарниры наклона:}
  двенадцать вершин $O_2(3)$ распадаются на три взаимно ортогональных экваториальных квадрата по $4$~вершины в координатных плоскостях $(xy,xz,yz)$, отвечающих трём углам Эйлера ($\dim\mathrm{SO}(3)=d(d-1)/2=3$), тогда как каждая из $8$~треугольных граней соединяет по одной вершине из плоскостей $xy$, $xz$ и $yz$, обеспечивая переход от горизонтального вращения к пространственному наклону за один шаг по ребру.
\item \textbf{Точная сферическая изотропия во втором порядке $O(|\mathbf{k}|^2)$ и точный закон анизотропии в четвёртом порядке $O(|\mathbf{k}|^4)$:}
  для любого волнового вектора $\mathbf{k}=(k_1,k_2,k_3)$ сумма степеней проекций по $12$ вершинам $e_j\in O_2(3)$ удовлетворяет точным тождествам
  \begin{equation}
    \label{eq:cuboct-moments-2-4}
    \sum_{j=1}^{12} (\mathbf{k}\cdot e_j)^2 \;=\; 8\,|\mathbf{k}|^2,\qquad
    \sum_{j=1}^{12} (\mathbf{k}\cdot e_j)^4 \;=\; 12\,|\mathbf{k}|^4 - 4\sum_{i=1}^3 k_i^4,
  \end{equation}
  вследствие чего дисперсионное разложение
  \begin{equation}
    \label{eq:cuboct-dispersion-2-4}
    \sum_{j=1}^{12}\bigl(1 - \cos(\mathbf{k}\cdot e_j)\bigr)
    \;=\; \underbrace{4\,|\mathbf{k}|^2}_{\text{круглая волна }O(|\mathbf{k}|^2)}
    \;-\; \underbrace{\left(\frac{1}{2}\,|\mathbf{k}|^4 - \frac{1}{6}\sum_{i=1}^3 k_i^4\right)}_{\text{кубическая анизотропия }O(|\mathbf{k}|^4)}
    \;+\; O(|\mathbf{k}|^6)
  \end{equation}
  строго изотропно во втором порядке, а решёточная анизотропия отодвинута в четвёртый порядок, где отношение $\sum_{i=1}^3 k_i^4/|\mathbf{k}|^4$ отличается ровно в $d=3$ раза между координатной осью $(1,0,0)$ ($12-4=8$) и пространственной диагональю $(1,1,1)$ ($12-4/3=32/3$, вдоль диагонали грани $(1,1,0)$ --- $12-2=10$).
\end{enumerate}
\end{proposition}

\begin{proof}
Двенадцать точек лежат на сфере радиуса~$a$ и попарно симметричны относительно координатных плоскостей.
Любая грань $\mathrm{conv}\,N(0)$ обязана лежать в плоскости, несущей максимальную аффинную оболочку подмножества~$N(0)$.
Такие плоскости --- координатные квадраты (четыре соседа в одной плоскости $x=\mathrm{const}$)
и равносторонние треугольники на диагональных сечениях $x+y+z=\mathrm{const}$.
Других вершин у выпуклой оболочки нет: все точки $N(0)$ уже вершины.
Полученное тело совпадает с классическим кубооктаэдром (дуальный ромбическому додекаэдру --- ячейке Вороного на \latGCKshort).
\end{proof}

\begin{remark}[Что не рассматривается]
Икосаэдр, додекаэдр или усечённый октаэдр \emph{не} возникают как $\mathrm{conv}\,N(x)$ у периодической решётки
с одной длиной ребра~$\caStepSpace$ и первой координационной сферой в~$\mathbb{R}^3$:
у Браве-носителя соседи --- фиксированный дискретный набор, а не произвольные двенадцать точек на сфере.
Ячейка Вороного на \latGCKshort --- ромбический додекаэдр; это \emph{зона одного узла}, а не однотактовая каузальная оболочка (\cref{def:one-tick-hull}).
\end{remark}

\subsection{Геометрия однотактового тела на \latGCKshort}
\label{sec:fcc-kappa}

Далее метрика тела $\mathcal{K}(x)$ из \cref{prop:cubocta-fcc}.
Вершины --- те же точки $(\pm a/\sqrt{2},\pm a/\sqrt{2},0)$ и перестановки, где $a=\caStepSpace$.
Расстояние от центра до вершины уже сосчитано:
\[
  R_{\mathrm{out}}=a.
\]
Четыре точки $(a/\sqrt{2},\pm a/\sqrt{2},0)$ и $(a/\sqrt{2},0,\pm a/\sqrt{2})$ лежат в плоскости $x=a/\sqrt{2}$ и образуют квадрат со стороной~$a$.
Таких квадратных граней шесть.
Расстояние от центра до грани
\[
  R_{\mathrm{in}}^{\mathrm{hull}}=a/\sqrt{2}.
\]
Отношение двух длин одного и того же тела
\[
  \caKappaGeom=R_{\mathrm{in}}^{\mathrm{hull}}/R_{\mathrm{out}}=1/\sqrt{2}
\]
ещё не использует скорость света: это геометрия оболочки.
Треугольная грань через $(a/\sqrt{2},a/\sqrt{2},0)$, $(a/\sqrt{2},0,a/\sqrt{2})$ и $(0,a/\sqrt{2},a/\sqrt{2})$ лежит в плоскости $x+y+z=a\sqrt{2}$ и отстоит от центра на $a\sqrt{2/3}$.
Квадратные и треугольные грани лежат на разных расстояниях, поэтому одной вписанной сферы, касающейся всех граней, нет.
В мост $\speedC=\caKappaGeom \caSignalSpeed$ входит расстояние до квадратной грани
(\cref{fig:carrier-cuboctahedron,fig:carrier-cubocta-meridian}).

\begin{figure}[htbp]
  \centering
  \input{figures/tikz/carrier-cuboctahedron.tex}
  \caption{Кубооктаэдр --- выпуклая оболочка двенадцати соседей \latGCKshort.}
  \label{fig:carrier-cuboctahedron}
\end{figure}

\begin{figure}[htbp]
  \centering
  \input{figures/tikz/carrier-cubocta-meridian.tex}
  \caption{Меридиан: квадратная и треугольная грани на разных расстояниях от центра.}
  \label{fig:carrier-cubocta-meridian}
\end{figure}

Объём ячейки Вороного считается отдельно.
Кубическая ячейка со стороной $A=a\sqrt{2}$ имеет объём $A^3=2\sqrt{2}\,a^3$ и содержит четыре узла \latGCKshort, поэтому на один узел
\[
  v_{\hv}=A^3/4=a^3/\sqrt{2}.
\]
Стена Вороного проходит посредине ребра к соседу, на расстоянии $a/2$ от узла.
Это не радиус оболочки:
\[
  R_{\mathrm{in}}^{\mathrm{Voronoi}}=a/2=\caKappaGeom\, R_{\mathrm{in}}^{\mathrm{hull}}.
\]
(\cref{fig:carrier-hull-voronoi,fig:carrier-voronoi-cell}).
Двенадцать соседей двойственны двенадцати ромбам этой ячейки.

\begin{figure}[htbp]
  \centering
  \input{figures/tikz/carrier-hull-voronoi.tex}
  \caption{Три радиуса однотактового тела на одном узле.}
  \label{fig:carrier-hull-voronoi}
\end{figure}

\begin{figure}[htbp]
  \centering
  \input{figures/tikz/carrier-voronoi-cell.tex}
  \caption{Замощение $\mathbb{R}^3$ ячейками Вороного (ромбододекаэдры).}
  \label{fig:carrier-voronoi-cell}
\end{figure}

Объём самой оболочки находится вычитанием.
Многогранник с вершинами всех перестановок $(\pm 1,\pm 1,0)$ есть пересечение октаэдра $|x|+|y|+|z|\le 2$ и куба $|x|,|y|,|z|\le 1$.
Объём октаэдра $|x|+|y|+|z|\le 1$ равен $4/3$, поэтому при правой части~$2$ объём равен $(4/3)\cdot 8=32/3$.
Шесть шапок за гранями куба, по одной на полуось, имеют объём $2/3$ каждая: при $x$ от~$1$ до~$2$ сечение $|y|+|z|\le 2-x$ есть квадрат в повёрнутых осях площади $2(2-x)^2$, и
\[
  \int_1^2 2(2-x)^2\,\mathrm{d}x=\frac{2}{3}.
\]
Шесть шапок дают~$4$, остаётся $32/3-4=20/3$.
Ребро этого многогранника равно $\sqrt{2}$.
Масштаб к ребру~$a$ есть множитель $a/\sqrt{2}$, объём умножается на $(a/\sqrt{2})^3$:
\[
  V_{\mathrm{hull}}=\frac{20}{3}\cdot\frac{a^3}{2\sqrt{2}}=\frac{5\sqrt{2}}{3}\,a^3,
  \qquad
  \frac{V_{\mathrm{hull}}}{v_{\hv}}=\frac{5\sqrt{2}}{3}\cdot\sqrt{2}=\frac{10}{3}.
\]
Формула $(8\sqrt{2}/3)\,a^3$ с отношением $16/3$ --- объём другого тела.

Макроскопическая скорость связана с тактовой тем же отношением радиусов оболочки:
\begin{align}
  c &= \caKappaGeom \caSignalSpeed, \\
  \htick &= \caKappaGeom\, t_P, \\
  \caSignalSpeed &= \sqrt{2}\,c,
  \label{eq:hull-speed-bridge}
\end{align}
где $t_P=\caStepSpace/c$ --- учебниковая метка.
Непрерывное определение $t_P=\caStepSpace/c$ предполагало бы $\caSignalSpeed=c$; на этой оболочке тактовая и макроскопическая скорости различаются множителем~$\sqrt{2}$.

